\section{Introduction}
\label{sec:intro}

Object detection~\cite{girshick2015fast, ren2016faster, redmon2016you, redmon2018yolov3, carion2020end, zhang2022dino, tian2019fcos, liu2016ssd, liu2022dab, zhao2024detrs, khanam2024yolov11, tan2020efficientdet, duan2019centernet} has long been a foundational task in computer vision due to its broad applications. The field has progressed from early CNN-based architectures, such as YOLO~\cite{redmon2016you} and Faster R-CNN~\cite{ren2016faster}, to Transformer-based models like DETR~\cite{carion2020end} and DINO~\cite{zhang2022dino}, while the task itself has evolved from traditional closed-set detection to open-set detection~\cite{liu2023grounding, li2022grounded, jiang2024t, cheng2024yolo, minderer2023scaling, ren2024dino, VLP:MDETR, minderer2022simple, minderer2023scaling} to better handle emerging real-world challenges.

A paramount objective in object detection is to develop models capable of identifying arbitrary objects and concepts. A common approach to this problem is open vocabulary object detection, where models such as Grounding DINO~\cite{liu2023grounding} and DINO-X~\cite{ren2024dino} leverage text encoders (e.g., BERT~\cite{devlin2018bert} or CLIP~\cite{radford2021clip}) to represent object categories and perform category-level open-set detection. Despite their effectiveness, these methods are fundamentally constrained by their relatively shallow language understanding, which restricts their ability to handle complex semantic descriptions (In Figure~\ref{fig:intro}, Grounding DINO detected all apples, despite the input prompt being \textit{red apple}.). Consequently, these methods are inherently limited in fully addressing this objective.

In contrast, multimodal large language models (MLLMs)~\cite{gpt4v,  VLM:Gemini, lu2025ovis2, wang2025internvl3, chen2022pali, agrawal2024pixtral, li2024aria, deitke2024molmo, wang2024qwen2} benefit from the strong language understanding capabilities of their underlying LLMs, presenting a promising avenue for integrating advanced language comprehension into object detection. A common MLLM-based approach~\cite{chen2023shikra, you2023ferret, wang2023cogvlm, zhan2025griffon, zhang2024ferret, jiang2024chatrex, ma2024groma, jiang2025referringperson, zhu2025internvl3, guo2025seed1, bai2025qwen2} is to represent coordinates as discrete tokens~\cite{chen2021pix2seq} and predict bounding boxes through next-token prediction. While conceptually elegant, existing MLLM-based methods have rarely matched the performance of traditional regression-based detectors on benchmarks such as COCO. As exemplified in Figure \ref{fig:intro}, even advanced MLLMs like Qwen2.5-VL~\cite{bai2025qwen2} struggles with precise object localization, in addition to facing limitations such as low recall rates, coordinate drift, and duplicate predictions.

\begin{figure}[t]     
    % \hspace*{-0.05\textwidth}
    \includegraphics[width=1.0\textwidth]{figures/intro.pdf}
    \caption{1) Detectors excel in localization but lack language understanding. MLLMs understand language well but struggle with localization. 2) Differences in optimization difficulty between detectors and MLLMs.}
    \label{fig:intro}
\end{figure}

We argue that the performance disparity in MLLM-based object detection primarily arises from two fundamental challenges inherent in their current formulation and training. First, MLLMs typically treat coordinate prediction as a discrete classification task, directly generating absolute coordinate values and relying on cross-entropy loss for supervision. While traditional regression-based models benefit from continuous, geometry-aware losses (e.g., L1, GIoU) that are directly sensitive to small geometric offsets, MLLMs face a significant learning difficulty in accurately mapping a fixed set of discrete tokens to the continuous pixel space. As illustrated in Figure \ref{fig:intro}, even minor pixel misalignments in discrete coordinate predictions can result in disproportionately large cross-entropy losses, hindering precise localization. This inherent challenge underscores the need for effective strategies to reduce coordinate learning complexity and provide extensive data for this mapping. 

Secondly, MLLMs commonly employ Supervised Fine-tuning (SFT) for teacher-guided next-token prediction training~\cite{radford2018improving}. While efficient, this paradigm creates a fundamental mismatch between training and inference. During SFT, the model is always conditioned on a ground-truth prefix, namely teacher forcing, meaning it is never exposed to its own, potentially imperfect, predictions. This training setup fails to capture the model's true performance in an autonomous generation setting.
This inherently prevents the model from developing robust behavioral awareness. Consequently, during free-form inference without this direct guidance, the model often struggles to regulate its own output structure. This leads to anomalous coordinate sequence generation, such as spurious duplicate predictions or object omissions, which undermine its overall performance. Addressing these two intertwined challenges is crucial for advancing MLLM-based object detection.



To overcome these inherent limitations and unlock the full potential of MLLMs for precise and versatile object perception, we propose \textbf{Rex-Omni}, a 3B-scale MLLM that achieves competitive performance with traditional detectors while distinctly excelling in language understanding capabilities. We address the aforementioned challenges through three core design principles:
\vspace{-1em}

\begin{itemize}
    \item \textbf{Task Formulation}: We unify visual perception tasks under a coordinate prediction framework, wherein each task is formulated as generating a sequence of coordinates. Specifically, pointing predicts one point, detection employs two points to form a bounding box, polygons use four or more points to represent object contours, and keypoint tasks output multiple semantic points. We adopt a quantized coordinate representation, where each coordinate value is mapped to one of 1,000 discrete tokens corresponding to values from 0 to 999. This approach significantly reduces the coordinate learning complexity and eases optimization, concurrently enhancing the efficiency of spatial representation.

    \item \textbf{Data Engines}: To facilitate the model's learning of the mapping between 1,000 discrete coordinate tokens and pixel-level positions, and to foster robust comprehension of complex natural language expressions, we design multiple specialized data engines for grounding, referring, and pointing tasks. These engines generate high-quality, semantically rich visual supervision signals for coordinate prediction.

    \item \textbf{Training Pipelines}:  We adopt a two-stage training paradigm. In the first stage, we perform supervised fine-tuning on 22 million data to teach the model basic coordinate prediction skills. In the second stage, we apply GRPO-based~\cite{deepseekmath} reinforcement post-training with three geometry-aware reward functions. This reinforcement phase serves two purposes: it enhances the precision of coordinate predictions through continuous geometric supervision, and crucially, it mitigates undesired behaviors (such as duplicate predictions) that arise from the teacher-guided nature of the initial SFT stage.
\end{itemize}

After this two-stage training, Rex-Omni achieves superior performance across a diverse range of perception tasks, as shown in Figure \ref{fig:teaser}., including object detection, object referring, visual prompting, GUI grounding, layout grounding, OCR, pointing, keypointing, and spatial referring. All of these tasks are achieved through direct prediction of coordinate points. To quantitatively assess its performance, Rex-Omni is first evaluated on COCO~\cite{lin2014microsoft}, a core benchmark for object detection. In a zero-shot setting (without training on COCO data), Rex-Omni demonstrates superior F1-score performance compared to traditional coordinate regression-based models (e.g., DINO-ResNet50, Grounding DINO) and other MLLMs (e.g., SEED1.5-VL~\cite{guo2025seed1}). Beyond COCO, Rex-Omni's performance is further benchmarked across diverse tasks, such as long-tailed detection, referring expression comprehension, dense object detection, GUI grounding, and OCR. Rex-Omni consistently outperforms both traditional detectors and MLLMs, thereby establishing a unified framework that combines precise localization with robust language understanding.

In summary, Rex-Omni represents a significant step towards unifying robust language understanding with precise visual perception. By carefully integrating principled task formulations, advanced data engines, and a sophisticated two-stage training pipeline, we demonstrate that MLLMs have the profound potential to define the next generation of object detection models, offering unprecedented versatility and a truly language-aware approach to visual perception systems.